\begin{proposition}[6.10.6]
\label{IV.2.s6.10.6-fr}
Sean $X$ un preesquema localmente noetheriano, $\mathcal F$ un
$\mathcal O_X$-Módulo coherente, $Y$ un subpreesquema cerrado irreducible de
$X$, con punto genérico $y\in\operatorname{Supp}(\mathcal F)$. Existe
entonces una vecindad abierta no vacía $U$ de $y$ en $X$ tal que, ara todo
$x\in U\cap Y$, se tiene
\begin{align}
\dim(\mathcal F_x)&=\dim(\mathcal F_y)+\dim(\mathcal O_{Y,x})
\tag{6.10.6.1}\label{IV.2.s6.10.6.1-fr}\\
\operatorname{prof}(\mathcal F_x)&=\operatorname{prof}(\mathcal F_y)+
\operatorname{prof}(\mathcal O_{Y,x})
\tag{6.10.6.2}\label{IV.2.s6.10.6.2-fr}\\
\operatorname{coprof}(\mathcal F_x)&=\operatorname{coprof}(\mathcal F_y)+
\operatorname{coprof}(\mathcal O_{Y,x}).
\tag{6.10.6.3}\label{IV.2.s6.10.6.3-fr}
\end{align}
\end{proposition}

\begin{proof}
\oldpage[IV-2]{158}
Sea $Y'=Y_{\mathrm{red}}$, que es un subpreesquema cerrado \emph{íntegro} de $Y$
con el mismo espacio subyacente, definido por un Ideal localmente nilpotente
$\mathcal K$ de $\mathcal O_Y$ (\textbf I, 6.1.6). Se sigue que
$\operatorname{gr}_{\mathcal K}^{\bullet}(\mathcal O_Y)$ es un
$\mathcal O_{Y'}$-Módulo \emph{coherente}, y, como $Y'$ es íntegro, existe una
vecindad abierta $V$ de $y$ en $Y'$ tal que
$\operatorname{gr}_{\mathcal K}^{\bullet}(\mathcal O_Y)|V$ es localmente
libre (\textbf{0}$_{\rm I}$, 5.2.7). Es decir, reemplazando $X$ por una
vecindad de $y$, se puede suponer que $\mathcal O_Y$ es \emph{normalmente plano}
a lo largo de $Y'$; se deduce, pues, de (6.10.3) que se tiene
\[
\dim(\mathcal O_{Y,x})=\dim(\mathcal O_{Y',x}),\qquad
\operatorname{prof}(\mathcal O_{Y,x})=
\operatorname{prof}(\mathcal O_{Y',x})
\]
para todo $x\in Y$; esto permite limitarse al caso en que el subpreesquema
cerrado $Y$ es \emph{íntegro}.

En estas condiciones, en virtud de (6.10.2), reemplazando $X$ por una
vecindad abierta de $y$, se puede suponer que $\mathcal F$ es \emph{normalmente plano} a lo
largo de $Y$. Como $y\in\operatorname{Supp}(\mathcal F)$, la relación
(6.10.6.1) se sigue de (6.10.5.2) y (6.10.5.4). Pongamos ahora
$p=\operatorname{prof}(\mathcal F_y)$; reemplazando, si es necesario, $X$ por una
vecindad abierta de $y$, se puede suponer que existen $p$ secciones $f_i$
$(1\leq i\leq p)$ de $\mathcal O_X$ sobre $X$, que forman una sucesión
$\mathcal F$-regular, tales que los $(f_i)_y$ pertenecen al ideal
maximal $\mathfrak m_y$ de $\mathcal O_{X,y}$ y que
$\operatorname{prof}(\mathcal F_y/(\sum_{i=1}^{p}(f_i)_y\mathcal F_y))=0$
(\textbf 0, 16.4.6). Si $\mathcal J$ es el Ideal de $\mathcal O_X$ que define
$Y$, se tiene $\mathcal J_y=\mathfrak m_y$ y, reemplazando de nuevo, si es necesario, $X$
por una vecindad de $y$, se puede suponer que
$f_i\in\Gamma(X,\mathcal J)$ para $1\leq i\leq p$. Además, si se pone
$\mathcal G=\mathcal F/(\sum_{i=1}^{p}f_i\mathcal F)$, la hipótesis
$\operatorname{prof}(\mathcal G_y)=0$ implica que $\mathcal G_y$ contiene
un submódulo isomorfo a
$\mathcal O_{X,y}/\mathfrak m_y$ (\textbf 0, 16.4.6); el mismo razonamiento (teniendo
en cuenta (\textbf{0}$_{\rm I}$, 5.3.8 y 5.2.7)) muestra que se puede suponer que
existe un \emph{sub-$\mathcal O_X$-Módulo} $\mathcal G'$ de $\mathcal G$
\emph{isomorfo a} $\mathcal O_X/\mathcal J$, de modo que
$\mathcal G'_x=\mathcal O_{Y,x}$ para todo $x\in Y$. Pongamos
$\mathcal G''=\mathcal G/\mathcal G'$; reemplazando $X$ por una vecindad
abierta de $y$, se puede, en virtud de (6.10.2), suponer que $\mathcal G'$ y
$\mathcal G''$ son \emph{normalmente planos} a lo largo de $Y$. Sea entonces $x$ un
punto cualquiera de $Y$, y pongamos
$q=\operatorname{prof}(\mathcal O_{Y,x})=\operatorname{prof}(\mathcal G'_x)$;
sea $(g_j)_{1\leq j\leq q}$ una sucesión maximal
$\mathcal G'_x$-regular de elementos del ideal maximal
$\mathfrak m_x$ de $\mathcal O_{X,x}$. Cada componente homogénea de
$\operatorname{gr}_{\mathcal J_x}^{\bullet}(\mathcal G''_x)$ es un
$\mathcal O_{Y,x}$-módulo plano finitamente generado por hipótesis y, por tanto, un
$\mathcal O_{Y,x}$-módulo \emph{libre}, y lo mismo ocurre, pues, con
$\operatorname{gr}_{\mathcal J_x}^{\bullet}(\mathcal G''_x)$; como la
sucesión $(g_j)$ es $\mathcal O_{Y,x}$-regular, es, por tanto,
$\operatorname{gr}_{\mathcal J_x}^{\bullet}(\mathcal G''_x)$-regular,
de donde se concluye que es $\mathcal G''_x$-\emph{regular} (\textbf 0, 15.1.19).
Aplicando el lema (\textbf 0, 15.1.18) a la sucesión exacta
\[
0\longrightarrow\mathcal G'_x\longrightarrow\mathcal G_x
\longrightarrow\mathcal G''_x\longrightarrow0
\]
por inducción sobre $j$, se obtiene una sucesión exacta
\[
0\longrightarrow\mathcal G'_x/
\left(\sum_{j=1}^{q}g_j\mathcal G'_x\right)
\longrightarrow\mathcal G_x/
\left(\sum_{j=1}^{q}g_j\mathcal G_x\right).
\]
Pero, por hipótesis,
$\operatorname{prof}(\mathcal G'_x/(\sum_{j=1}^{q}g_j\mathcal G'_x))=0$
(\textbf 0, 16.4.6); por la caracterización (\textbf 0, 16.4.6) de los módulos de profundidad
nula, se concluye que
$\operatorname{prof}(\mathcal G_x/(\sum_{j=1}^{q}g_j\mathcal G_x))=0=\operatorname{prof}(\mathcal F_x/(\sum_{i=1}^{p}f_i\mathcal F_x+\sum_{j=1}^{q}g_j\mathcal F_x))$;
como la sucesión $(g_j)$ es $\mathcal G'_x$-regular y
$\mathcal G''_x$-regular, es también $\mathcal G_x$-\emph{regular}; finalmente, la
sucesión $((f_i)_x)$ es $\mathcal F_x$-regular por hipótesis y está formada
por elementos del ideal maximal $\mathfrak m_x$; se deduce (\textbf 0, 16.4.6)
que $\operatorname{prof}(\mathcal F_x)=p+q$, lo que concluye la
demostración.
\end{proof}
